\def\loge{Log(E)} 
\chapter{Introduction}\label{Introduction}
Spin models are, in the most general sense, systems in which variables with two or more possible states can have interactions with the rest of the system. This description may be overly reductive, but does reflect how versatile spin models can be. Spin models can be used to model magnetism, such as with the Ising model, or to mimic cell behavior, such as in cellular Potts models.

The two most fundamental discriminators between different spin models are how many possible states a variable can have and its spin operators, or how the variables in the system interact. This thesis will use binary datasets and focus on how the variable interactions. For example, the Ising model can be seen as a pairwise binary model, as interactions occur on only one or two variables at a time. 

The flexibility of spin models allows for the construction of arbitrarily complex models, but this may not necessarily be desirable. As \citet{bayesianinference} have pointed out, an arbitrarily complex model may be able to describe any dataset, but they lack the explanatory power of simpler models, as these models may be easier to study and falsify.\@ \citet{complexity} have studied the complexity and found that the simplest models focus their parameters on a small number of variables. 

With this knowledge, the concept of the Minimally Complex Model (MCM) has been developed.~\cite{bayesianinference} In an MCM, variables are divided into groups called Independent Complete Components. As the name suggests, these ICCs have no interactions between each other, but are complete internally. This means that every interaction between states of any order is included in the model, from interactions of only a single spin to one which contains all variables in the ICC.\@ By quantifying the complexity of a given MCM for a dataset, an algorithmic search can then be performed to find which MCM fits the dataset best. 

The complexity can be quantified by calculating the Bayesian evidence, or marginal likelihood, for the combination of an MCM and the dataset. The evidence of an MCM is given by
\begin{equation}
  P(\hat{s}|\mathcal{M})= \frac{1}{2^{N(n-r)}} \prod_{a \in \mathcal{A}} \left[ \frac{\Gamma(2^{r_{a} - 1})}{\Gamma(N+2^{r_{a} -1})} \prod_{b_{a} } \left[ \frac{\Gamma(k_{b_{a} }+\frac{1}{2})}{\Gamma\left( \frac{1}{2} \right)} \right] \right]
\end{equation}
which consists of three parts. The first part compensates for the rank of the MCM, with \(N\) the amount of datapoints in the dataset \(\hat{s}\), \(n\) the amount of variables in each datapoint, and \(r\) the rank of the MCM, which is the amount of variables of the dataset included in the model; usually \(n = r\), making the starting term one. The first product then iterates over every ICC in the MCM and calculates a compensation term for the rank of each ICC \(r_a\). Finally, the second product iterates over each of the \(2^{r_a}\) possible datapoints the ICC could see, and uses the amount of them in the dataset \(k_{b_a}\). This fraction also becomes one when \(k_{b_a}=0\), so only datapoints actually present in the dataset contribute here.

Because the calculated probabilities here get very small, the natural logarithm of probability is taken, which gives the natural logarithm of the evidence, also called the \loge{} of the model. Because \(P(\hat{s}|\mathcal{M})<1\), the \loge{} is a negative number, where a higher (closer to zero) number indicates a lower complexity. Taking the logarithm also has the additional benefit of turning the products and division in the equation into sums and subtractions, which are more efficient to compute. 
